The only incoming theorem is the entire qualified G331, with its unchanged paper, annotation, ledger and scope. Let \(R=A_{\mathrm{red}}\) for excellent local \(A\), \(n=\dim R\ge1\), and let \(\mathfrak a\) be proper primary for the marked maximal ideal \(\mathfrak m\). Write \(X=X_{\mathfrak a}\) for the full compact Hausdorff section \(v(\mathfrak a)=1\). Every prime support is retained. G331 supplies the independently geometric continuous class \(C\), its equality with the uniform ideal closure, and its actual decreasing-net completion \(D\). All functions here are nonpositive; products take bounded \(D\) inputs. G331 grants no operator.

All multiplicities use degree \(n\) of the REDUCED ring. Modules of smaller dimension contribute zero. A positive-dimensional local ring has no zero-dimensional minimal component: its maximal ideal cannot be minimal while a smaller prime exists. Thus the zero-dimensional minimal-component exception in G331's valuative criterion does not arise here. Superficial quotients may be nonreduced, nonpure, and have embedded primes.

The ordinary published multigraded Hilbert theorem (complete current July 2024 Swanson--Huneke 17.4.2 and Definition 17.4.3) says that for finite \(M\) and proper primary \(I_1,\ldots,I_r\), the length \(\ell_R(M/I_1^{s_1}\cdots I_r^{s_r}M)\) is eventually a rational polynomial of degree \(\dim M\). Its degree-\(n\) form is
\[
H_M(s)=\sum_{|d|=n}\frac{e_R(I_1^{[d_1]},\ldots,I_r^{[d_r]};M)}{d_1!\cdots d_r!}s^d.
\]
With \(n\) separately labelled entries, \(e(I_1,\ldots,I_n)\) is the coefficient of \(s_1\cdots s_n\). Repeated equal entries recover ordinary multiplicity by substituting the sum of their variables. Unit labels mean independence of their variables and zero coefficients whenever those variables occur.

Unused labels need proof. Fix the unused primary \(B\)'s exponent at a sufficiently large \(c\). If \(K\) is a remaining primary label, choose \(N\) with \(K^N\subset B^c\). For the remaining product \(P_s\),
\[
P_sK^N\subset B^cP_s\subset P_s.
\]
The outside length polynomials differ only by a constant exponent shift and have the same top form. Squeezing on every positive rational ray proves the same top form for the middle polynomial. In the full polynomial, fixing the unused exponent leaves a degree-\(n\) part in the other variables coming exactly from terms with unused exponent zero. Polynomial identity proves deletion of unused labels, including repeated labels. With no remaining labels the degree-\(n>0\) coefficient is zero; units are immediate.

Relabelling proves symmetry. Substitution \(s\mapsto ls\) proves \(e(I^l,K_2,\ldots,K_n)=l e(I,K_2,\ldots,K_n)\). In the polynomial with \(n+1\) labels \(I,J,K_2,\ldots,K_n\), set the first two variables equal to \(s\). Only exponent vectors \((1,0,1,\ldots,1)\) and \((0,1,1,\ldots,1)\) contribute to \(s s_2\cdots s_n\). The unused-label result identifies their coefficients, giving
\[
e(IJ,K_2,\ldots,K_n)=e(I,K_2,\ldots,K_n)+e(J,K_2,\ldots,K_n).
\]
Iteration proves every power/product identity, with unit factors.

For integral-closure invariance suppose \(I\subset J\) and \(J\) is integral over \(I\). Homogeneous ideal equations make \(R[Jt]\) finite over \(R[It]\). Homogeneous module generators of degrees at most \(c\) give \(J^{m+c}\subset I^m\): in degree \(m+c\) their coefficients lie in \(I^{m+c-d}\subset I^m\). For \(m>c\) and any other primary product \(P_s\),
\[
I^mP_s\subset J^mP_s\subset I^{m-c}P_s.
\]
The outside length polynomials have identical top forms. Their difference on positive rational rays is lower order, so the middle form agrees on every such ray, hence identically. All mixed coefficients agree. Equal integral closures reduce to this case twice through \(I+J\). This is equality of forms, not a coefficient inequality deduced from a nonnegative polynomial; \((s-t)^2\) demonstrates why the latter inference fails.

For all-field monotonicity suppose \(J\subset I\) and keep \(K_2,\ldots,K_n\) fixed. Pass algebraically, if needed, to \(T=R[z]_{\mathfrak mR[z]}\), faithfully flat local with residue field \(k(z)\). Tensoring a finite-length composition series preserves length, with factors \(k(z)\). Dimension is preserved: an \(n\)-generator parameter ideal extends to an \(\mathfrak mT\)-primary ideal, giving \(\dim T\le n\) by height, while flat going-down lifts an \(n\)-chain below \(\mathfrak mT\), giving the reverse bound. Length polynomials and coefficients are therefore unchanged. No boundary base-change claim is made.

Over the infinite residue field, complete 17.2.2 provides complements of finitely many proper linear subspaces in \(K_n/\mathfrak mK_n\) for superficial choices for each of \((K_n,J,K_2,\ldots)\) and \((K_n,I,K_2,\ldots)\). Add avoidance of the current minimal primes. These added subspaces are proper: \(K_n\subset P+\mathfrak mK_n\) would imply \(K_n\subset P\) by Nakayama on \(K_n/(K_n\cap P)\), impossible for a primary ideal in positive dimension. Their finite intersection has a point over an infinite field. Choose one lift \(x\), using the SAME quotient for both families. Its dimension is \(d-1\): a quotient prime chain extends strictly down to a minimal prime avoiding \(x\), giving the upper bound; lifting a quotient parameter system with \(x\) gives the lower bound by height.

The colon calculation controls embedded primes. Multi-ideal Artin--Rees follows by finite homogeneous generation of the intersection module over the multi-Rees algebra. For superficial \(x\), its consequence (complete 17.2.4), using the common primary radical and Krull intersection, is
\[
(P_s:x)=(0:x)+P_{s-e_n},\qquad (0:x)\cap P_{s-e_n}=0
\]
for all large exponents. Artin--Rees first puts the colon in \((0:x)+P_{s-c}\); a sufficiently high \(K_n\)-power puts the remainder in the truncated range where superficiality applies. Intersecting the truncated superficial identities over arbitrarily large \(K_n\)-exponents and applying Krull intersection proves the second assertion. Thus \((0:x)\) injects into a finite primary quotient and has finite length. Multiplication by \(x\) gives
\[
\ell(R/P_s)-\ell(R/P_{s-e_n})
=\ell(R/(x,P_s))-\ell(0:x).
\]
For \(d>1\) the final constant does not affect the degree-\((d-1)\) coefficients. Reading the coefficient with one \(K_n\) entry removed gives the exact superficial reduction for BOTH families on the same quotient (complete 17.4.6). The unused-label result removes the zero \(K_n\) label.

Repeat in \(K_{n-1},\ldots,K_2\), always using simultaneous choices in the current common quotient. All images stay proper primary, and the arbitrary Noetherian quotient hypotheses remain valid. The final one-dimensional ring \(R'\) gives \(e_{R'}(JR')\) and \(e_{R'}(IR')\). The inequality between lengths of their powers, divided by the exponent, proves
\[
e_R(J,K_2,\ldots,K_n)\ge e_R(I,K_2,\ldots,K_n).
\]
For \(n=1\) this is direct. The same reduction and positive one-dimensional Samuel leading coefficient prove positive mixed multiplicities for proper entries. Units give zero. Symmetry proves every-slot monotonicity, and integral-closure invariance gives order for included integral closures. No regular-sequence or Cohen--Macaulay assumption was inserted.

We also need exact-sequence top-form additivity. For \(0\to M_1\to M_2\to M_3\to0\), multi-ideal Artin--Rees sandwiches \(P_sM_2\cap M_1\) between \(P_sM_1\) and \(P_{s-c}M_1\). The ray squeeze proves the same top form for its quotient lengths. The quotient exact sequence gives \(H_{M_2}=H_{M_1}+H_{M_3}\), with lower-dimensional modules contributing zero. This proves the application of complete 17.4.4 without asserting an injection after reduction modulo \(P_s\).

Finally G331's primary criterion on ALL \(X\) says \(v(J)\ge v(K)\) iff \(J\subset\overline K\), with units handled separately. Equal rational profiles give equal integral closures of the corresponding cross-powers. The laws above give equal mixed moments; profile inequalities give reversed ideal-moment order. These are mathematical proofs with the stated ordinary published Hilbert and superficial dependencies, not finite experiments.
